Fix \(n\in\mathbb N\), \(1\leq\beta\leq n\), \(B>0\), and
\(\omega=(\omega_1,\ldots,\omega_\beta)\in(0,\infty)^\beta\).  The inequalities
\(1\leq\beta\leq n\) imply \(n\geq1\).  By
Definition~\(\mathrm{def{:}ellipsoidal\text{-}coefficient\text{-}class}\), membership of
\((Y,c)\) in \(\mathcal C_{n,\beta}(B,\omega)\) includes membership of
\((Y,c)\) in \(\mathcal C^{\mathrm{CEY}}_{n,\beta}\).  Therefore
\[
\mathcal C_{n,\beta}(B,\omega)
\subseteq \mathcal C^{\mathrm{CEY}}_{n,\beta}.
\tag{1}
\]

We prove strictness by the witness in the proposition.  Since \(n\geq1\) and
\(\beta\geq1\), the singleton \(\{1\}\) belongs to
\(\mathcal S_{n,\beta}\).  For every \(i\in I_n\), set
\[
c^*_{i,\{1\}}=2B\sqrt{\frac{\omega_1}{n}},
\qquad
c^*_{i,R}=0
\quad
\text{for all }R\in\{\varnothing\}\cup\mathcal S_{n,\beta}
\text{ with }R\neq\{1\},
\tag{2}
\]
and define, for every \(z\in\{0,1\}^n\),
\[
Y_i^*(z)
=c^*_{i,\varnothing}
+\sum_{R\in\mathcal S_{n,\beta}}c^*_{i,R}
  \prod_{j\in R}z_j
=2B\sqrt{\frac{\omega_1}{n}}\,z_1.
\tag{3}
\]
The square root and the division in (2)--(3) are well defined because
\(n\geq1\) and \(\omega_1>0\).  Equation~(3) verifies
Assumption~\(\mathrm{ass{:}low\text{-}order\text{-}schedule}\) for the fixed array
\((Y^*,c^*)\).  Hence Definition~\(\mathrm{def{:}published\text{-}low\text{-}order\text{-}class}\)
gives
\[
(Y^*,c^*)\in\mathcal C^{\mathrm{CEY}}_{n,\beta}.
\tag{4}
\]

For each \(R\in\mathcal S_{n,\beta}\), the definition
\(a_R=n^{-1}\sum_{i\in I_n}c_{i,R}\), together with
\(|I_n|=n\) and (2), yields
\[
a_R^*
=\frac1n\sum_{i\in I_n}c^*_{i,R}
=
\begin{cases}
2B\sqrt{\omega_1/n},&R=\{1\},\\
0,&R\neq\{1\}.
\end{cases}
\tag{5}
\]
By the definition of \(\Omega\), this matrix is diagonal with entry
\(\omega_{|R|}>0\) in coordinate \(R\).  Thus positivity of every order weight
implies that \(\Omega\) has full rank and
\((\Omega^{-1})_{R,R}=\omega_{|R|}^{-1}\).  Substituting (5) gives
\[
(a^*)^{\mathsf T}\Omega^{-1}a^*
=\sum_{R\in\mathcal S_{n,\beta}}
  \frac{(a_R^*)^2}{\omega_{|R|}}
=\frac{\bigl(2B\sqrt{\omega_1/n}\bigr)^2}{\omega_1}
=\frac{4B^2}{n}.
\tag{6}
\]
Because \(B>0\) and \(n\geq1\),
\[
\frac{4B^2}{n}>\frac{B^2}{n}.
\tag{7}
\]
Equations~(6)--(7) show that the witness violates the ellipsoid condition
\(a^{\mathsf T}\Omega^{-1}a\leq B^2/n\) from
Assumption~\(\mathrm{ass{:}aggregate\text{-}coefficient\text{-}ellipsoid}\) and
Definition~\(\mathrm{def{:}ellipsoidal\text{-}coefficient\text{-}class}\).  Consequently,
\[
(Y^*,c^*)\notin\mathcal C_{n,\beta}(B,\omega).
\tag{8}
\]
Combining (1), (4), and (8) proves
\[
\mathcal C_{n,\beta}(B,\omega)
\subsetneq\mathcal C^{\mathrm{CEY}}_{n,\beta},
\]
with exactly the strictness witness asserted in the proposition.